This study constructed a reproducible normative indicator of field-of-study mismatch among employed STEM university graduates using annual ENEMDU data from 2021--2025. Of the 24,525 evaluable person-period records, the weighted mismatch estimate for 2025 was 40.28\%. Differences by field and sex, together with lower median nominal income and less favourable employment conditions in the mismatched group, demonstrate the value of examining occupational compatibility alongside general employment indicators.

The machine learning models showed modest predictive performance. Logistic regression achieved the highest ROC-AUC, XGBoost the highest AP and Random Forest the highest balanced accuracy; no clear overall winner was established. Sex, STEM field and province contributed consistently to prediction. The sensitivity analyses support the stability of several findings under changes in training specification, but do not resolve the limitations of the normative label or longitudinal linkage. These findings provide a basis for further investigation of educational and occupational alignment rather than for individual employment decisions or causal policy recommendations.
